\documentclass[11pt,a4paper]{article}
\usepackage[utf8]{inputenc}
\usepackage[T1]{fontenc}
\usepackage[english]{babel}
\usepackage{geometry}
\usepackage{graphicx}
\usepackage{listings}
\usepackage{xcolor}
\usepackage{hyperref}
\usepackage{amsmath}
\usepackage{booktabs}
\usepackage{enumitem}

\geometry{margin=2.5cm}
\hypersetup{colorlinks=true, linkcolor=violet, urlcolor=violet}

\definecolor{bilpurple}{HTML}{4A2D7A}
\definecolor{codebg}{HTML}{F5F0FA}

\lstset{
  backgroundcolor=\color{codebg},
  basicstyle=\ttfamily\small,
  breaklines=true,
  frame=single,
  rulecolor=\color{bilpurple},
}

\title{
  \textcolor{bilpurple}{\textbf{Term Sheet Extraction Pipeline}}\\
  \large Architecture, Algorithms, and Technical Documentation
}
\author{Nasser Kassioui}
\date{March 2026}

\begin{document}
\maketitle
\tableofcontents
\newpage

% ==========================================================================
\section{Introduction}

This document describes the architecture and algorithms of a coordinate-based
extraction pipeline for structured product term sheets in PDF format. The system
processes documents in English, French, and German, and is optimised for
BIL (Banque Internationale \`a Luxembourg) term sheets.

The core insight driving the architecture is that \textbf{the problem is not
text extraction, but structure reconstruction}. Standard PDF text extraction
tools destroy the spatial relationships between labels and values. Our pipeline
recovers this structure by working directly with word coordinates.

% ==========================================================================
\section{Architecture Overview}

The pipeline consists of six modules:

\begin{enumerate}[label=\arabic*.]
  \item \textbf{coord\_extractor} --- Extracts words with bounding boxes,
        reconstructs lines and columns, builds label-value pairs.
  \item \textbf{lang\_detector} --- Detects document language (EN/FR/DE)
        using keyword frequency.
  \item \textbf{field\_extractor} --- Two-pass field extraction: structured
        pairs first, regex fallback second.
  \item \textbf{excel\_exporter} --- Excel output with BIL purple styling.
  \item \textbf{highlighter} --- Annotates PDFs with coloured highlights.
  \item \textbf{app.py} --- Streamlit web interface.
\end{enumerate}

Data flows linearly: PDF $\rightarrow$ coord\_extractor $\rightarrow$
lang\_detector $\rightarrow$ field\_extractor $\rightarrow$ excel\_exporter.

% ==========================================================================
\section{Coordinate-Based Extraction}

\subsection{The Word Gluing Problem}

\texttt{pdfplumber}'s default \texttt{x\_tolerance=3} merges words whose
characters are within 3 points of each other. BIL term sheets use tight
character spacing, producing outputs like:

\begin{lstlisting}
SettlementCurrency  (instead of: Settlement Currency)
ISINCH1481964653    (instead of: ISIN CH1481964653)
\end{lstlisting}

\textbf{Solution}: We set \texttt{x\_tolerance=1, y\_tolerance=1}. This
produces 745 properly separated words on page 1 of the EN test document,
versus 107 glued blobs with defaults.

\subsection{Line Reconstruction}

Words are sorted by (page, top, x0) and grouped into lines using a
Y-proximity threshold:

$$
\text{same\_line}(w_i, w_j) \iff |w_i.\text{top} - w_j.\text{top}| \leq \tau_y
$$

where $\tau_y = 3.0$ points. This value was determined empirically from the
BIL document layout where line height is 7--12pt and inter-line gap is 4--8pt.

\subsection{Column Detection}

For each page, we build a histogram of word left-edge positions ($x_0$),
binned to 5-point resolution. The two most frequent bins separated by at
least 20 points become column boundaries.

In BIL term sheets, this reliably detects:
\begin{itemize}
  \item Column 1 (labels): $x_0 \approx 36$ pt
  \item Column 2 (values): $x_0 \approx 142$ pt
\end{itemize}

\subsection{Label-Value Pair Construction}

For each reconstructed line, words left of the column boundary form the
\textit{label}; words right form the \textit{value}. Lines with labels
exceeding 60 characters are classified as paragraph text and skipped.

% ==========================================================================
\section{Field Extraction}

\subsection{Two-Pass Strategy}

\textbf{Pass 1} (structured): Each label-value pair is looked up in a
multilingual label map (\texttt{config.LABEL\_TO\_FIELD}) that covers
English, French, and German variants.

\textbf{Pass 2} (regex fallback): If Pass 1 does not find a field, regex
patterns are applied to the full reconstructed text.

This ordering matters: structured extraction from coordinates is more
reliable than regex on linearised text, because it preserves the spatial
relationship between labels and values.

\subsection{ISIN Classification}

All valid ISINs are found via the pattern \texttt{[A-Z]\{2\}[A-Z0-9]\{10\}}.
The product's own ISIN (PST\_ISIN) is distinguished from underlying ISINs
by proximity scoring: the ISIN closest to an ``ISIN'' label in the text
is classified as PST\_ISIN.

\subsection{Multilingual Support}

Each field has regex patterns covering all three languages. For example,
the Issuer field matches:
\begin{itemize}
  \item EN: ``Issuer''
  \item FR: ``\'Emetteur'' / ``Emetteur''
  \item DE: ``Emittentin''
\end{itemize}

% ==========================================================================
\section{Difficulties Encountered}

\begin{enumerate}
  \item \textbf{Word gluing} was the single biggest obstacle. Discovering
        that \texttt{x\_tolerance=1} solves it was the breakthrough.
  \item \textbf{Multi-word labels split across columns}: labels like
        ``Conditional Coupon Amount'' span the column boundary, placing
        ``Conditional'' in the label column and ``Coupon Amount 5.60\%''
        in the value column. We handle this by searching both label and
        value for field-specific keywords.
  \item \textbf{German declensions}: ``Schlechteste'' (nominative) vs
        ``Schlechtesten'' (genitive) required wildcard suffixes in regex.
  \item \textbf{Absent fields}: Barrier Reverse Convertibles (SSPA 1230)
        do not specify Capital Protection. The system correctly returns
        None rather than guessing.
\end{enumerate}

% ==========================================================================
\section{Validation Results}

Tested on 3 BIL term sheets (EN, FR, DE):

\begin{table}[h]
\centering
\begin{tabular}{lcc}
\toprule
\textbf{Field} & \textbf{Fill Rate} & \textbf{Accuracy} \\
\midrule
PST\_ISIN         & 100\% & 100\% \\
BIL               & 100\% & 100\% \\
ISSUER            & 100\% & 100\% \\
CURRENCY          & 100\% & 100\% \\
MATURITY          & 100\% & 100\% \\
COUPON            & 100\% & 100\% \\
WORST\_OR\_AVERAGE & 100\% & 100\% \\
SSPA\_TYPE        & 100\% & 100\% \\
\bottomrule
\end{tabular}
\caption{Extraction results on the 3-document test set.}
\end{table}

\end{document}
